\chapter{Préférences et disposition}
\section{Ouvrir les préférences}
Utiliser le menu \emph{Préférences} ou \touche{Ctrl + virgule}. Les trois onglets
sont \emph{Interface}, \emph{Apparence} et \emph{Avancé · RAW}.
Les préférences d’interface sont distinctes des retouches de chaque photo.

\section{Interface}
Afficher ou masquer Informations, Histogramme et Ruban ; choisir la largeur
du panneau Réglages (260--360 px) et la taille du texte des informations
(10--20 px). Les mêmes panneaux peuvent être masqués depuis le menu Affichage.

Les poignées entre les panneaux latéraux et la vue image permettent de
redimensionner Informations (210--300 px) et Réglages. Celle au-dessus du ruban
ajuste sa hauteur (100--260 px). Glisser à la souris, ou donner le focus à
la poignée et utiliser les flèches au pas de 10 px. Un double-clic restaure
la taille usuelle. Les dimensions sont mémorisées dans l’INI.

\section{Apparence}
Les palettes proposées sont Bordeaux, Blanc cassé, Noir mat, Bleu marine et
Gris actuel. Elles s’appliquent aux panneaux et aux préférences. Le fond de
la photo est indépendant : changer de palette ne remplace pas ce choix.
Les fonds bleu marine et rouge bordeaux utilisent des dégradés.

\attention{Pour juger finement la couleur d’une photo, un fond gris neutre
limite l’influence de l’entourage coloré. La palette de l’interface ne change
pas les pixels exportés et ne constitue pas un profil écran ICC.}

\section{Avancé · RAW}
L’onglet donne accès à la calibration de l’image ouverte, au manifeste et à
la réinitialisation profil/calibration. Ce dernier bouton est distinct de
\emph{Réinitialiser les réglages}, qui remet les retouches à leurs références
mais n’efface pas automatiquement tous les choix de développement RAW.
La réinitialisation générale est également accessible avec
\touche{Ctrl + Maj + R} et demande confirmation avant d’agir sur la photo ouverte.

\section{Persistance et état des réglages}
Les sections dépliées, dimensions, couleurs et certains choix d’interface
sont mémorisés. Le chemin du fichier INI apparaît dans les préférences.
Pour une édition manuelle, fermer IPR et conserver une copie du fichier.

Un nom de contrôle numérique en gras et un point sur sa section indiquent
une différence par rapport à sa valeur de réinitialisation. Ce repère ne
couvre pas encore toutes les courbes, cases et recettes ; ce n’est pas
un indicateur exhaustif de fichier enregistré ou non enregistré.
Les mémos explicatifs au survol prolongé restent prévus pour la phase bêta.
